\subsection{The LEGB Rule Invariant}

Unqualified name lookup follows Local $\rightarrow$ Enclosing $\rightarrow$ Global $\rightarrow$ Built-in. The compiler does not implement this as a runtime dictionary walk for every access inside functions. It classifies names at compile time and emits specialized bytecodes.

\subsubsection{Variable Scope and Lexical Boundaries: The LEGB Resolution Rules}

For names bound in the current function scope, CPython emits \texttt{LOAD\_FAST} / \texttt{STORE\_FAST}, indexing a fixed array of local slots in the frame. Access is $O(1)$ array indexing with no hash lookup.

For globals and builtins, the compiler emits \texttt{LOAD\_GLOBAL} / \texttt{STORE\_GLOBAL} (or \texttt{LOAD\_NAME} at module level), which perform dictionary lookups in \texttt{f\_globals} and the builtins namespace. These paths carry hash-table overhead on every access.

%<<<PROGRAM:programs/the01>>>


Enclosing-scope reads from nested functions use \texttt{LOAD\_DEREF} / \texttt{STORE\_DEREF} on cell objects---heap-allocated indirection slots shared between outer and inner code objects.

Shadowing is a compile-time classification consequence: if a name is assigned anywhere in a function body, the compiler treats it as local for the entire function, so reads before assignment raise \texttt{UnboundLocalError} rather than falling through to global lookup.

Built-in names (\texttt{len}, \texttt{print}) resolve through \texttt{LOAD\_GLOBAL} unless shadowed by a local or global binding. The \texttt{builtins} module holds the outermost namespace; \texttt{import builtins; builtins.len(...)} bypasses shadowing.

The LEGB invariant governs which namespace wins; the bytecode instruction (\texttt{LOAD\_FAST} vs \texttt{LOAD\_GLOBAL} vs \texttt{LOAD\_DEREF}) reveals the compiler's chosen fast or slow path for each name.
